% !TEX program = xelatex
% 공통수학2 · Ⅱ. 집합과 명제 · 2. 명제 · 03 명제의 증명 · 1/2차시
\documentclass[11pt,a4paper]{article}
\usepackage{kotex}
\usepackage{iftex}
\ifPDFTeX\else
  \setmainhangulfont{Noto Serif CJK KR}
  \setsanshangulfont{Noto Sans CJK KR}
\fi
\usepackage[margin=2.1cm]{geometry}
\usepackage{parskip}
\usepackage{enumitem}
\usepackage{array}
\usepackage{tabularx}
\usepackage{booktabs}
\usepackage{xcolor}
\usepackage{amssymb}
\usepackage{tikz}
\usepackage[most]{tcolorbox}
\usepackage{fancyhdr}
\usepackage{titlesec}

\definecolor{goldc}{HTML}{B07C22}
\definecolor{goldstrong}{HTML}{8A5F16}
\definecolor{indigoc}{HTML}{2B3B57}
\definecolor{indigosoft}{HTML}{6E82A6}
\definecolor{linec}{HTML}{D6DACB}
\definecolor{surface2}{HTML}{E4E8DC}
\definecolor{notebg}{HTML}{FBF3E3}
\definecolor{inksoft}{HTML}{52604F}

\titleformat{\section}{\normalfont\sffamily\bfseries\large\color{indigoc}}{\thesection}{0.6em}{}
\titlespacing{\section}{0pt}{16pt}{8pt}
\newtcolorbox{ideabox}{enhanced, breakable, colback=white, colframe=goldc, boxrule=0.5pt, leftrule=3pt, arc=2pt, left=10pt, right=10pt, top=8pt, bottom=8pt}
\newtcolorbox{calloutbox}{enhanced, breakable, colback=white, colframe=indigosoft, boxrule=0.5pt, leftrule=3pt, arc=2pt, left=10pt, right=10pt, top=7pt, bottom=7pt}
\newtcolorbox{notebox}{enhanced, breakable, colback=notebg, colframe=goldc, boxrule=0.5pt, leftrule=3pt, arc=2pt, left=10pt, right=10pt, top=7pt, bottom=7pt}
\newtcolorbox{answerbox}{enhanced, breakable, colback=white, colframe=linec, boxrule=0.5pt, arc=2pt, left=10pt, right=10pt, top=7pt, bottom=7pt}
\newcommand{\lessonbanner}[3]{%
  \begin{tcolorbox}[enhanced, colback=indigoc, colframe=indigoc, boxrule=0pt, arc=0pt,
    left=14pt, right=14pt, top=14pt, bottom=10pt, borderline south={2.4pt}{0pt}{goldc}, fontupper=\color{white}]
    {\sffamily\footnotesize\color{white!85!black} #1}\\[4pt]{\sffamily\bfseries\Large #2}\\[3pt]{\sffamily\small\color{white!88!black} #3}
  \end{tcolorbox}}
\pagestyle{fancy}\fancyhf{}\renewcommand{\headrulewidth}{0pt}
\fancyfoot[L]{\footnotesize\color{inksoft} 공통수학2 · 2026학년도 2학기 · 지평선고등학교}
\fancyfoot[R]{\footnotesize\color{inksoft} 교사용 지도안 -- 배포 금지}

\begin{document}
\lessonbanner{공통수학2 · Ⅱ. 집합과 명제 · 2. 명제}{03 명제의 증명 --- 1/2차시}{대우를 이용한 증명법과 귀류법 · 교사용 지도안 (2026학년도 2학기)}
\vspace{10pt}

\noindent\begin{tabularx}{\linewidth}{@{}>{\bfseries\color{indigoc}}p{2.6cm}X@{}}
\toprule
대상 & 고1 · 2개 반 · 반당 13명 \\
차시 & 50분 (도입 10 · 전개 30 · 정리 10) \\
성취기준 & 대우를 이용한 증명법과 귀류법을 이해하고, 관련된 명제를 증명할 수 있다. \\
선행 학습 & 29차시 --- 충분조건과 필요조건 \\
\bottomrule
\end{tabularx}

\section{학습 목표 \& Big Idea}
\begin{ideabox}
\textbf{\color{goldstrong}영속적 이해 (Big Idea)}\\
명제 $p\to q$가 참임을 직접 보이기 어려울 때, 논리적으로 동치인 다른 명제 --- 대우, 또는 `결론을 부정하면 모순이 생긴다'는 사실 --- 를 보이면 간접적으로 증명할 수 있다.
\vspace{4pt}
\small\color{inksoft} 핵심개념: \textbf{간접증명} \quad\textbullet\quad 핵심개념: \textbf{모순} \quad\textbullet\quad 일반화: \textbf{동치인 명제를 증명해도 원래 명제가 증명된다}
\end{ideabox}
\begin{calloutbox}
\textbf{차시 학습 목표}
\begin{itemize}[leftmargin=1.4em, itemsep=2pt, topsep=4pt]
  \item 대우를 이용하여 명제가 참임을 증명할 수 있다.
  \item 귀류법을 이용하여 명제가 참임을 증명할 수 있다.
\end{itemize}
\end{calloutbox}

\section{질문의 연결}
\begin{tabularx}{\linewidth}{@{}>{\bfseries\color{indigoc}\small}p{2.4cm}X@{}}
사실적 질문 & 대우를 이용한 증명법과 귀류법은 각각 어떤 절차로 진행될까? \\[6pt]
개념적 질문 & 왜 결론을 부정하는 두 방법(대우 증명, 귀류법)이 서로 다른 것을 보여 주는데도 둘 다 원래 명제를 참으로 만들 수 있을까? \\[6pt]
논쟁적 질문 & 직접 보일 수 없는 것을 `모순이 생긴다'는 사실만으로 참이라고 결론짓는 귀류법은 정당한 추론일까? \\
\end{tabularx}

\section{수업 흐름}
\begin{tabularx}{\linewidth}{@{}>{\centering\arraybackslash}p{1.6cm}X>{\raggedright\arraybackslash\small}p{3.1cm}@{}}
\toprule
\textbf{단계} & \textbf{활동 \& 발문} & \textbf{ATL / VTR} \\
\midrule
\textcolor{goldstrong}{\textbf{도입}}\newline\footnotesize 10분 &
\textbf{마음공부 (2분)} --- 새 소단원 시작, 유무념 대조.\newline
\textbf{생각 열기 (8분)} --- ``지역 주민 또는 교복을 입은 사람에게 입장료를 할인해 준다.'' $\to$ ``교복을 입은 사람은 할인받는다.''(참) $\to$ 그 대우 ``할인받지 못한 사람은 교복을 입고 있지 않다.''도 참인지 판단. &
ATL: 사고(논리 전이)\newline VTR: See-Think-Wonder \\
\addlinespace
\textcolor{goldstrong}{\textbf{전개}}\newline\footnotesize 30분 &
\textbf{대우를 이용한 증명법 (15분)} --- 명제와 대우의 참·거짓이 일치함을 복습 $\to$ 예제 1(`$n^2$이 짝수이면 $n$도 짝수'를 대우로 증명) $\to$ 문제 1(`$n^2$이 3의 배수이면 $n$도 3의 배수', $n=3k-2$·$n=3k-1$ 경우 나누기).\newline
\textbf{귀류법 (15분)} --- 결론(또는 명제)의 부정을 가정해 모순을 이끄는 방법 정의 $\to$ 예제 2($\sqrt2$는 무리수) $\to$ 문제 2($1+\sqrt2$는 무리수). &
ATT: 촉진자 역할\newline ATL: 논리적 유도 \\
\addlinespace
\textcolor{goldstrong}{\textbf{정리}}\newline\footnotesize 10분 &
\textbf{개념 연결 질문 (5분)} --- ``대우를 이용한 증명과 귀류법은 둘 다 `부정'에서 시작하는데, 무엇이 다를까?'' (대우 증명: 결론의 부정 $\to$ 가정의 부정을 보임 / 귀류법: 결론의 부정 $\to$ 이미 아는 사실과의 모순을 보임)\newline
\textbf{성찰 (5분)} --- 감사 일기 1문장 + 다음 차시 예고(절대부등식). &
마음공부 · 감사 일기 \\
\bottomrule
\end{tabularx}

\begin{calloutbox}
\textbf{지도상의 유의점} --- 대우를 이용한 증명과 귀류법은 둘 다 `결론을 부정한다'는 점에서 비슷해 보이지만, 대우 증명은 결론을 부정하면 가정이 부정됨을 보이는 것이고, 귀류법은 결론을 부정하면 가정이나 이미 참인 사실과 모순됨을 보이는 것임을 구별하여 지도한다. 직접증명법(정의·공리·정리로 가정에서 결론을 직접 이끎)과 간접증명법(대우·귀류법)의 차이도 짚어 준다.
\end{calloutbox}

\section{형성평가 \& 답안}
\begin{answerbox}
\textbf{문제 1} \quad `자연수 $n$에 대하여 $n^2$이 3의 배수이면 $n$도 3의 배수이다.' --- 대우로 증명\\
\textbf{\color{goldstrong}답} \; 대우: $n$이 3의 배수가 아니면 $n^2$도 3의 배수가 아니다. $n$이 3의 배수가 아니면 $n=3k-2$ 또는 $n=3k-1$($k$는 자연수)로 나타낼 수 있다.\\
\small (ⅰ) $n=3k-2$일 때 $n^2=9k^2-12k+4=3(3k^2-4k+1)+1$이므로 $n^2$도 3의 배수가 아니다.\\
(ⅱ) $n=3k-1$일 때 $n^2=9k^2-6k+1=3(3k^2-2k)+1$이므로 $n^2$도 3의 배수가 아니다.\\
따라서 대우가 참이므로 주어진 명제도 참이다.
\end{answerbox}

\vspace{6pt}
\begin{answerbox}
\textbf{문제 2} \quad `$1+\sqrt2$는 유리수가 아니다.' --- 귀류법으로 증명\\
\textbf{\color{goldstrong}답} \; 명제를 부정하여 $1+\sqrt2$가 유리수라 하면 $1+\sqrt2=a$($a$는 유리수)로 나타낼 수 있다. 즉 $\sqrt2=a-1$이고, 유리수끼리의 뺄셈은 유리수이므로 $a-1$은 유리수이다. 그런데 좌변의 $\sqrt2$는 무리수이므로 모순이다. 따라서 $1+\sqrt2$는 유리수가 아니다.
\end{answerbox}

\section{성취수준 (이 차시 범위)}
\begin{tabularx}{\linewidth}{@{}>{\bfseries\color{goldstrong}}p{0.9cm}X@{}}
\toprule
A & 대우를 이용한 증명법과 귀류법을 이해하고 관련된 명제를 스스로 증명할 수 있다. \\
B & 대우를 이용한 증명법과 귀류법을 알고 관련된 명제의 증명을 부분적으로 수행할 수 있다. \\
C & 대우를 이용한 증명법과 귀류법을 알고 안내된 절차에 따라 부분적으로 증명할 수 있다. \\
D & 안내된 절차에 따라 대우를 이용한 증명을 부분적으로 수행할 수 있다. \\
E & 대우를 이용한 증명법의 의미를 안다. \\
\bottomrule
\end{tabularx}

\section{다음 차시}
\begin{calloutbox}
\textbf{03 명제의 증명 --- 2/2차시.} 절대부등식의 뜻을 알고, 산술평균과 기하평균의 관계 및 절댓값을 포함한 절대부등식을 증명하는 활동을 다룬다.
\end{calloutbox}
\end{document}
